% answers_short.tex —— 2005 年·普通高等学校招生考试·江西卷·理科 · 简短答案速查
% 由 tools/make_short_answers.py 自 answers.tex 的 \gksolhead/\gkanswer 提取，
% 并与 2005/2005_数学真题_江西理科_答案要点.md 逐题比对后生成，勿手改（改源 MD 后重跑）。

\begin{gknotice}
\textbf{使用说明}：本节只给最终答案便于快速核对，完整推导见解析版。答案为编写者独立推导，\textbf{非官方标准答案}。
\end{gknotice}

\gkpart{选择题}
\begin{center}\begin{tabular}{ccccccccccccc}
\toprule
题号 & 1 & 2 & 3 & 4 & 5 & 6 & 7 & 8 & 9 & 10 & 11 & 12 \\
\midrule
答案 & D & A & D & B & B & C & C & C & C & B & D & A \\
\bottomrule
\end{tabular}\end{center}

\gkpart{填空题}
\begin{itemize}[leftmargin=2.6em]
  \item[\textbf{13.}] $\dfrac{\sqrt{2}}{2}$。
  \item[\textbf{14.}] $\dfrac{3}{2}$。
  \item[\textbf{15.}] $\dfrac{3\sqrt{2}}{2}$。
  \item[\textbf{16.}] ③④。
\end{itemize}

\gkpart{解答题}
\begin{itemize}[leftmargin=2.6em]
  \item[\textbf{17.}] (1) $f(x) = \dfrac{x^2}{2-x}$；(2) 当 $1 < k < 2$ 时解集为 $(1, k) \cup (2, +\infty)$；当 $k = 2$ 时解集为 $(1, 2) \cup (2, +\infty)$；当 $k > 2$ 时解集为 $(1, 2) \cup (k, +\infty)$
  \item[\textbf{18.}] 不存在这样的实数 $x \in [0, \pi]$。
  \item[\textbf{19.}] (1) $\xi$ 的取值范围为 $\{5, 7, 9\}$；(2) $\mathbf{E}\xi = \dfrac{275}{32}$
  \item[\textbf{20.}] (1) 证明见解析，由 $A_1D \perp AD_1$ 且 $A_1D \perp AB$ 证得 $A_1D \perp \text{\text{平面} } AD_1B$；(2) $\dfrac{1}{3}$；(3) $2 - \sqrt{3}$
  \item[\textbf{21.}] (1) 证明见解析，由数学归纳法结合函数单调性证得；(2) $a_n = 2\left[1 - \left(\dfrac{1}{2}\right)^{2^n}\right]$（或 $a_n = 2 - \left(\dfrac{1}{2}\right)^{2^n - 1}$）
  \item[\textbf{22.}] (1) $y = \dfrac{4}{3}x^2 - \dfrac{1}{3}x + \dfrac{2}{3}$（或 $3y = 4x^2 - x + 2$）；(2) 证明见解析，由抛物线切线性质或数量积等式 $\vec{FA}\cdot\vec{FP} = |FA|(y_0 + \frac{1}{4})$ 证得
\end{itemize}
